\chapter{The Four Operations of Distinction}
\label{ch:four-operations}

\begin{quotation}
\noindent\textit{Synopsis.} Chapter~\ref{ch:distinctions} built a self-contained
apparatus for renderings and discard sets without situating it within the
broader research program from which its vocabulary descends. This chapter
does that situating, against six sources at three different levels of
verification: Distinguishability Geometry and Persistence Before Truth,
checked by reading primary text directly; Observability Is Restorability,
likewise checked directly, which also surfaces an inconsistency between
two of the other source texts; Admissibility Before Truth and The
Geometry of Correction, integrated on the strength of formulas and
theorem statements transcribed directly from source PDFs by the author
rather than independently checked against primary source text; and The
Geometry of Witnesses, reported only from the author's prose summary and
cited at a correspondingly weaker tier still. Every claim is marked \imported{},
\transcribed{}, or \secondary{} according to which of these three tiers
actually supports it, and the distinction is maintained rather than
allowed to blur after the first citation. The honest finding is that
this monograph built working analogues of three of Distinguishability
Geometry's four primitive operations without recognizing them as such,
resolved a fourth case only a later source could characterize, and
discovered -- by the accident of checking an unrelated source -- that
its oldest working distinction, benign versus corrupting loss, already
had a name in a program it had not previously engaged.
\end{quotation}

\section{Two Collisions in the Source Vocabulary}
\label{sec:collisions}

Before any importing can be done responsibly, two terminological collisions
need to be named, because silently picking one convention over the other
would misrepresent both source texts.

The first collision is internal to the wider research program this monograph
draws on. Distinguishability Geometry's foundational object is the ontological
triple $(X, \sim, T)$, in which $T$ denotes an \emph{ontology}: a description
language or template hierarchy determining which descriptions of elements of
$X$ are reachable. Persistence Before Truth's foundational object is the
structured domain $(X, \mathcal{T}, \mathcal{D}, \mathcal{R})$, in which
$\mathcal{T}$ denotes a family of \emph{admissible transformations}
$\{T_\alpha : X \to X\}$. The same letter is doing incompatible work in two
sibling foundational texts: one uses $T$ for the language in which $X$ is
described, the other uses $\mathcal{T}$ for the maps that act on $X$. This
monograph adopts the merged notation
\[
(X, \sim, T, \mathcal{T}),
\]
with $X$ the state space, $\sim$ the distinguishability relation, $T$ the
ontology in Distinguishability Geometry's sense, and $\mathcal{T}$ the
transformation family in Persistence Before Truth's sense, and uses this
notation for the remainder of the present chapter whenever both roles need to
appear together.

The second collision is not resolved here, and it is important to say so
plainly rather than let the first resolution create an impression of general
tidiness. This monograph's own admissible continuation set
$\admis{\tau}{r}$ (\cref{def:admissible-set}) is an epistemic, consumer-relative
notion: the continuations reachable given what a particular rendering
discloses. The Admissibility Program's own central object, so far as it can
be reconstructed from a downstream application paper and a secondary summary
rather than from the foundational text itself, is $\mathcal{A} \subseteq X$:
a structural, non-epistemic subset of states whose continuation would not
destroy the conditions for coherent further continuation, regardless of what
any particular consumer can currently tell. These are not the same kind of
object, and this monograph has not verified against the primary Admissibility
Field source what the precise relationship between them is. This is recorded
as an open item in Appendix~\ref{app:open-questions} rather than papered over.

\section{Two Deficits}
\label{sec:two-deficits}

\imported{Distinguishability Geometry}

\begin{definition}[Ontological Triple]
\label{def:ontological-triple}
An ontological triple is $(X, \sim, T)$, where $X$ is a set, $\sim$ is an
equivalence relation on $X$ (observational indistinguishability), and $T$ is
an ontology determining which descriptions of elements of $X$ are reachable.
The equivalence classes $[x]_\sim$ are called fibers.
\end{definition}

\begin{definition}[Coding Deficit]
\label{def:coding-deficit}
Let $L_T(x)$ be the shortest description of $x$ reachable within $T$, and
$L^*(x) = K(x)$ the Kolmogorov complexity of $x$. The coding deficit is
$\delta_T(x) = L_T(x) - L^*(x)$.
\end{definition}

\begin{definition}[Distinguishability Deficit]
\label{def:distinguishability-deficit}
Let $D_T(x)$ be the number of objects separable from $x$ by descriptions
reachable in $T$, and $D^*(x)$ the maximum such capacity attainable by any
ontology. The distinguishability deficit is $\Delta_T(x) = D^*(x) - D_T(x)$.
\end{definition}

\begin{theorem}[Deficit Proxy]
\label{thm:deficit-proxy}
\imported{Distinguishability Geometry}
Under faithful encoding (each expressible distinction requires at least one
description bit), $\Delta_T(x) \le C\,\delta_T(x)$ for a system-dependent
constant $C > 0$. In particular, $\delta_T(x) = 0 \Rightarrow \Delta_T(x) = 0$.
\end{theorem}

The converse direction is stated in the source text as an open conjecture,
not a theorem, and this monograph inherits that status rather than
strengthening it: whether $\Delta_T(x) = 0 \Rightarrow \delta_T(x) = 0$ in
general remains unproved.

This apparatus is considerably more developed than the loss vocabulary this
monograph built independently in Chapter~\ref{ch:distinctions}. Where
$\discard{r}$ (\cref{def:discard-set}) records only which named distinctions
are present or absent, $\delta_T$ and $\Delta_T$ are quantitative, and the
Deficit Proxy Theorem relates two genuinely different notions of loss --
coding cost and distinguishing capacity -- that this monograph never
separated. \cref{sec:where-this-sits} states precisely how far the
correspondence between the two vocabularies extends and where it does not.

\section{The Four Operations}
\label{sec:four-operations}

\imported{Distinguishability Geometry}

Distinguishability Geometry organizes representational change into two
levels. Projection and embedding operate at the distinction level, modifying
$\sim$ within a fixed space. Revision and transport operate at the ontology
level: revision changes $T$, and transport connects two distinct triples.

\begin{definition}[Projection]
\label{def:dg-projection}
A projection on $(X, \sim, T)$ is a surjective map $\sigma: X \to S$ inducing
a coarsening $\sim'$ of $\sim$: $\sigma(x) = \sigma(y) \Rightarrow x \sim' y$.
\end{definition}

\begin{theorem}[Projection Deficit Bound]
\label{thm:projection-deficit-bound}
\imported{Distinguishability Geometry}
Let $\sigma: X \to S$ be a projection inducing fiber partition $\{F_s\}_{s \in
S}$, with $X$ carrying a probability measure $\mu$. Then $\Delta\delta(\sigma)
\ge H(X \mid S)$, the conditional entropy of $X$ given $S$.
\end{theorem}

\begin{definition}[Embedding]
\label{def:dg-embedding}
An embedding of $(X, \sim, T)$ into $(Y, \sim', T')$ is an injective map
$\phi: X \to Y$ such that $\phi(x) \sim' \phi(y) \Rightarrow x \sim y$. An
embedding is faithful if the converse also holds.
\end{definition}

\begin{theorem}[Embedding Refinement]
\label{thm:embedding-refinement}
\imported{Distinguishability Geometry}
If $\phi: (X,\sim,T) \to (Y,\sim',T')$ is a faithful embedding, then
$\delta_{T'}(\phi(x)) \le \delta_T(x)$: a faithful embedding never increases
the coding deficit.
\end{theorem}

\begin{definition}[Revision]
\label{def:dg-revision}
An ontological revision is a transition $(X, \sim_T, T) \to (X, \sim_{T'},
T')$ where $T' \supset T$. A revision is productive if $\delta_{T'}(x) <
\delta_T(x)$.
\end{definition}

\begin{theorem}[Revision Monotonicity]
\label{thm:revision-monotonicity}
\imported{Distinguishability Geometry}
For an expanding sequence of ontologies $T_0 \subset T_1 \subset \cdots$,
$\delta_{T_{n+1}}(x) \le \delta_{T_n}(x)$ for all $n \ge 0$.
\end{theorem}

\begin{corollary}[Asymptotic Deficit]
\label{cor:asymptotic-deficit}
\imported{Distinguishability Geometry}
The sequence $(\delta_{T_n}(x))_{n \ge 0}$ is monotone non-increasing and
bounded below by zero; hence $\lim_{n\to\infty} \delta_{T_n}(x) = \delta_\infty(x)
\ge 0$ exists, with $\delta_\infty(x) = 0$ iff some reachable revision
sequence converges to an optimal ontology for $x$.
\end{corollary}

\begin{definition}[Transport]
\label{def:dg-transport}
A transport map is $\tau: (X,\sim_X,T_X) \to (Y,\sim_Y,T_Y)$ with $x \sim_X
x' \Rightarrow \tau(x) \sim_Y \tau(x')$, and distortion $\epsilon(\tau) \ge 0$
controlling the measure of fiber violations.
\end{definition}

\begin{theorem}[Transport Stability]
\label{thm:transport-stability}
\imported{Distinguishability Geometry}
Let $\tau$ be a transport map with distortion at most $\epsilon$, and let
$I_X$, $I_Y$ be invariants constant on fibers of $\sim_X$, $\sim_Y$
respectively, agreeing on perfectly transported fibers. Then $|I_X(x) -
I_Y(\tau(x))| \le C\epsilon$ for a constant $C$ depending on the oscillation
of the invariants across fiber boundaries.
\end{theorem}

\begin{corollary}[Justification of Analogy]
\label{cor:justification-of-analogy}
\imported{Distinguishability Geometry}
An analogical inference transported by $\tau$ is valid up to error
$C\epsilon(\tau)$; Transport Stability is the precise condition under which
analogical reasoning is valid.
\end{corollary}

\section{Persistence Sets}
\label{sec:persistence-sets}

\imported{Persistence Before Truth}

\begin{definition}[Structured Domain]
\label{def:structured-domain}
A structured domain is a quadruple $(X, \mathcal{T}, \mathcal{D}, \mathcal{R})$,
where $X$ is a state space, $\mathcal{T} = \{T_\alpha : X \to X\}$ is a
family of admissible transformations, $\mathcal{D} = \{d_i : X \to Y_i\}$ is
a family of distinctions, and $\mathcal{R} = \{R_j : X \to X\}$ is a family
of admissible reconstruction operators.
\end{definition}

\begin{definition}[Recoverability Relation]
\label{def:recoverability-relation}
A distinction $d \in \mathcal{D}$ is recoverable under $T \in \mathcal{T}$ if
there exists $R \in \mathcal{R}$ such that $d \circ R \circ T = d$.
\end{definition}

\begin{definition}[Persistence Set]
\label{def:persistence-set}
The persistence set of a distinction $d$ is $P(d) = \{T \in \mathcal{T} :
\exists R \in \mathcal{R} \text{ with } d \circ R \circ T = d\}$.
\end{definition}

\begin{theorem}[Persistence Depends Only on Structure]
\label{thm:persistence-structural}
\imported{Persistence Before Truth}
Recoverability is invariant under any isomorphism of structured domains that
preserves transformations, distinctions, and reconstruction operators; hence
recoverable persistence is a structural property rather than a property of
particular representations.
\end{theorem}

This is the general, primitive object this monograph's root-object apparatus
(\cref{def:root-object}) turns out to be a narrow special case of.
\cref{sec:where-this-sits} makes the correspondence precise.

\section{A Third Verified Source: Observability Is Restorability}
\label{sec:oir}

\imported{Observability Is Restorability}

A third source cited in this monograph's front matter has since been checked
against primary text, and the check surfaced something worth reporting
precisely: not an error this monograph introduced, but an inconsistency
already present between two source texts it draws on.

\begin{definition}[Restorable Distinction Set]
\label{def:rest-oir}
For a repair repertoire $\mathcal{R}$, $\mathrm{Rest}(\mathcal{R}) = \{d \in
\mathcal{D} : C_R(d) < +\infty\}$: the set of distinctions restorable at
finite cost under every admissible perturbation.
\end{definition}

$\restdomain$, used throughout Chapters~\ref{ch:repair-economics} onward, is
this object, and the attribution to Observability Is Restorability was
correct: $\mathrm{Rest}(\mathcal{R})$ is exactly this source's central
construction, arrived at from repair paths, perturbations, and a
substrate $(X_S, \mathcal{P}, \mathcal{D})$ considerably richer than
anything this monograph builds on top of it. This monograph's $\repcost{d}$
is a bare scalar; Observability Is Restorability's $C_R(d)$ is an infimum
over actual repair paths through a topological state space, and the source
text develops a full geometry on top of it -- an interference matrix with
spectral structure (bottleneck clusters, synergistic bundles), a repair
pseudometric, an observer-relative admissibility manifold $\mathcal{A}_O$
characterized by a Repair Boundary Theorem, and a Control-Geometry
Collapse theorem identifying restorability with reachability under
controllability hypotheses. None of this deeper apparatus is used in the
present monograph; $\restdomain$ is imported as a set, not as the geometry
that generates it.

\begin{remark}[An Inconsistency Between Two Source Texts, Not Introduced Here]
\label{rem:interference-inconsistency}
\emph{Conservation of Ambiguity} states that its repair interference
functional $I(d_i,d_j)$ is ``established in Observability Is
Restorability,'' and gives it in derivative form: the change in repair
cost of $d_j$ induced by an infinitesimal increase in repair investment
devoted to $d_i$. This is exactly the form \cref{def:interference} uses,
correctly attributed to \emph{Conservation of Ambiguity}. But Observability
Is Restorability's own Definition 8.7 gives a different, algebraic form:
$I(d_i,d_j) = C_R(d_i \wedge d_j) - C_R(d_i) - C_R(d_j)$, a superadditivity
gap rather than a derivative. The sign conventions agree (positive:
competition or conflict; negative: synergy; zero: independence), but the
formulas are not the same object, and one text's claim to inherit the
other's definition does not survive a direct check. This is recorded here
because it bears on how much weight the phrase ``established in
Observability Is Restorability'' can carry elsewhere in the corpus, not
because it requires this monograph to choose between the two: everything
in Chapters~\ref{ch:repair-economics} onward uses only the derivative
form, consistently, and nothing here depends on reconciling it with the
algebraic one.
\end{remark}

\begin{theorem}[Paradigm Shift as Repair-Repertoire Discontinuity]
\label{thm:paradigm-shift}
\imported{Observability Is Restorability}
Let $\mathcal{R}_{\mathrm{old}}$ and $\mathcal{R}_{\mathrm{new}}$ be repair
repertoires with $\mathcal{R}_{\mathrm{old}} \not\subseteq
\mathcal{R}_{\mathrm{new}}$ and $\mathcal{R}_{\mathrm{new}} \not\subseteq
\mathcal{R}_{\mathrm{old}}$. Then $\mathrm{Rest}(\mathcal{R}_{\mathrm{old}})
\not\subseteq \mathrm{Rest}(\mathcal{R}_{\mathrm{new}})$ and conversely: the
restorable-distinction sets are incomparable, with lost distinctions
$\mathrm{Rest}(\mathcal{R}_{\mathrm{old}}) \setminus
\mathrm{Rest}(\mathcal{R}_{\mathrm{new}})$ and gained distinctions the
reverse difference, both generically nonempty.
\end{theorem}

This theorem is the resolution \cref{prop:relay-replacement-partial-revision}
needed and did not have. That proposition showed relay revision matches
Distinguishability Geometry's revision operation only when the state space
is held fixed and the ontology merely expands, and left the genuine
paradigm-shift case -- where the state space itself changes --
uncharacterized. \cref{thm:paradigm-shift} characterizes exactly that case,
and does so within the repair-theoretic vocabulary this monograph already
uses throughout Chapters~\ref{ch:repair-economics} onward, rather than
requiring an appeal to Distinguishability Geometry's ontology-revision
apparatus at all. \cref{sec:where-this-sits} restates this as a fifth
correspondence.

\section{A Fourth Source, Transcribed: Admissibility Before Truth}
\label{sec:abt}

\transcribed{Admissibility Before Truth}

The following is imported on the strength of exact formulas and a proof
sketch transcribed directly from the source PDF by the author, not from
an independent reading of the primary file; per the front matter, this is
marked throughout rather than allowed to blend into the fuller confidence
of \cref{sec:oir}'s verified import.

\begin{definition}[Domain Restriction]
\label{def:domain-restriction}
A truth predicate $\mathcal{T}$ factors as $\mathcal{T} = \widetilde{\mathcal{T}}
\circ \iota$, where $\iota: \mathrm{Rest}(\mathcal{R}) \hookrightarrow
\mathcal{D}$ is the inclusion of restorable distinctions into the full
distinction space.
\end{definition}

\begin{theorem}[Domain Restriction Theorem]
\label{thm:domain-restriction}
\transcribed{Admissibility Before Truth}
$\mathrm{dom}(\mathcal{T}) \subseteq \mathrm{Rest}(\mathcal{R})$: a truth
predicate is defined only on distinctions restorable by the observer's
repair repertoire. Outside $\mathrm{Rest}(\mathcal{R})$, a distinction is
neither true nor false but unevaluable.
\end{theorem}

\begin{proofsketch}
\transcribed{Admissibility Before Truth} Every distinction in the domain
of $\mathcal{T}$ must lie in $\mathcal{D}$ by well-typedness; $\mathcal{T}$
is then defined on $\mathrm{Rest}(\mathcal{R})$ by restriction through
$\iota$, and outside $\mathrm{Rest}(\mathcal{R})$ the predicate takes a
third value marking non-evaluability, not a truth value of either
polarity.
\end{proofsketch}

\begin{proposition}[Sufficiency Presupposes Evaluability]
\label{prop:sufficiency-presupposes-evaluability}
Task-relative sufficiency (\cref{def:sufficiency}) implicitly assumes
$\mathcal{T}$'s domain restriction: $r$ can be sufficient for $\tau$ only
if every distinction $\tau$ requires already lies in
$\mathrm{Rest}(\mathcal{R})$ for the repair repertoire generating $r$.
\end{proposition}

\begin{proofsketch}
\cref{def:sufficiency} asks whether $\admis{\tau}{r} \subseteq
\mathcal{C}_\tau(e)$, which presupposes that membership in
$\mathcal{C}_\tau(e)$ is a meaningful, evaluable question for the
distinctions at stake. By \cref{thm:domain-restriction}, that question is
only meaningful for distinctions in $\mathrm{Rest}(\mathcal{R})$; for a
distinction outside it, asking whether a continuation is correct is
type-confused rather than merely undetermined. This was never stated in
Chapter~\ref{ch:admissibility}, which built \cref{def:sufficiency} without
reference to any prior evaluability condition, silently assuming it. The
assumption is almost always harmless in this monograph's worked domains,
since the distinctions at stake (castling rights, causal history, lexical
content) are restorable by construction whenever they are discussed at
all -- but it is an assumption, made visible only now.
\end{proofsketch}

\cref{prop:sufficiency-presupposes-evaluability} is the closest this
monograph comes to making contact with the Admissibility Program's actual
formal content, as distinct from the still entirely unseen
\emph{Admissibility Field}. It does not resolve the three-way
admissibility collision recorded in the front matter and in
Appendix~\ref{app:open-questions}, but it is the first of the four
competing admissibility-adjacent notions now on record that connects to
this monograph's own apparatus by a stated proof rather than by analogy
alone.

\section{A Fifth Source, Transcribed: The Geometry of Correction}
\label{sec:goc}

\transcribed{The Geometry of Correction}

\begin{definition}[Self-Referential and Reality-Anchored Evaluation]
\label{def:evaluation-types}
For an evaluation function $E: X \to \mathbb{R}$ and internal signal set
$S$, a system is self-referential if $E(x) = f(S(x))$, depending only on
internal signals. A system is reality-anchored if $E(x) = f(S(x), O(x))$
for an external channel $O$ the system does not fully author or control.
\end{definition}

\begin{definition}[Coupling Coefficient]
\label{def:coupling-coefficient}
$\rho = \partial E / \partial O$. When $\rho \approx 0$, the external
channel no longer moves the evaluation in practice, regardless of whether
$O$ remains present.
\end{definition}

\begin{principle}[Correction Suppression]
\label{principle:correction-suppression}
\transcribed{The Geometry of Correction}
Pathological self-reference does not usually mean no external corrective
channel exists; it means the system has become insensitive to a channel
that remains present.
\end{principle}

This is stated as a Principle rather than a Theorem because the source
text itself names it that way, and this monograph preserves that label
rather than upgrading its apparent formality.

\begin{proposition}[Corruption as Correction Suppression]
\label{prop:corruption-as-suppression}
Corrupting loss (\cref{def:corrupting-loss}) is an instance of correction
suppression: a rendering's presentation obscuring a discard is exactly a
case where the rendering's apparent evaluation of its own sufficiency has
decoupled from the external channel -- the actual, undisclosed state of
$\discard{r}$ -- that would otherwise correct it.
\end{proposition}

\begin{proofsketch}
Let $E$ be the confidence or sufficiency a rendering presents to a
consumer, and let $O$ be the true discard set $\discard{r}$, a fact about
the rendering the consumer does not directly observe. Declared loss
(\cref{def:benign-loss}) keeps $E$ coupled to $O$: the consumer's
assessment of sufficiency tracks what was actually discarded, since it is
disclosed. Undeclared loss (\cref{def:corrupting-loss}) is exactly the
case where $E$ no longer moves with $O$: the presented confidence is
computed from $S$ alone -- the rendering's internal, self-referential
account of itself -- with $\rho \approx 0$ relative to the actual discard
set. \cref{prop:confidence-renderings}'s regress (Chapter~\ref{ch:fanout})
is the recursive version of the same failure: each confidence profile
$Q$ is itself a rendering whose own $E$ may or may not track its own
$O$, and the trust boundary (\cref{def:trust-boundary}) is precisely the
point at which this monograph stops asking whether $\rho \approx 0$ has
crept in one level further down.
\end{proofsketch}

\cref{prop:corruption-as-suppression} is offered as the tightest
correspondence in this chapter, tighter even than the four/five-operations
correspondences of \cref{sec:where-this-sits}: it did not require
reinterpreting either definition, only recognizing that
\cref{def:corrupting-loss}'s ``obscures'' and \cref{def:coupling-coefficient}'s
``$\rho \approx 0$'' are the same phenomenon named twice, once in this
monograph's discard-set vocabulary and once in the correction-theoretic
vocabulary of a program this monograph had not previously engaged.

\section{A Sixth Source, at a Weaker Tier: The Geometry of Witnesses}
\label{sec:witnesses}

\secondary{The Geometry of Witnesses}

What follows rests on the author's prose account of having read this
text, not on transcribed formulas or primary source text, and is
reported at that correspondingly weaker tier: informative enough to
record and connect to this monograph's open questions, not solid enough
to state as a Definition or Theorem the way \S\S\ref{sec:abt}--\ref{sec:goc}
do.

The reported central question is who maintains the repair repertoire
$\mathcal{R}$ that Admissibility Before Truth (\S\ref{sec:abt}) treats as
given. The proposed answer introduces a witness operator, reconstructing
a repair operator from historical structure rather than repairing a
distinction directly, and a reconstruction cover $C(r)$ -- the set of
histories from which a given repair operator can be reconstructed -- with
an associated redundancy, reported as $\rho(r) = |C(r)|$.

\begin{remark}[A Notation Collision, Not Introduced Here]
\label{rem:rho-collision}
The Geometry of Witnesses's reported redundancy and The Geometry of
Correction's coupling coefficient (\cref{def:coupling-coefficient}) both
use $\rho$ for entirely unrelated quantities, apparently a collision
between the two source texts' own notational choices rather than an
error introduced in this monograph. To avoid compounding the collision
inside a single chapter, this monograph writes the witness redundancy as
$\rho_W(r) = |C(r)|$ from this point on, reserving unadorned $\rho$ for
\cref{def:coupling-coefficient}'s coupling coefficient throughout.
\end{remark}

Zero, $\rho_W(r) = 0$, means the repair operator is permanently lost; one
means a single point of failure; larger values mean resilience.

\begin{remark}[A Candidate Answer to an Open Question, Unverified]
\label{rem:witnesses-persistence-measure}
Appendix~\ref{app:open-questions} records, as unresolved, what
$\mu(P(d))$ -- the measure of a persistence set (\cref{def:persistence-set})
-- might be in general. The reconstruction cover $C(r)$ and redundancy
$\rho_W(r)$, as reported, are a plausible candidate operationalization of
exactly this quantity for the specific case of root recovery
(\cref{prop:root-recovery}): a root object with many independent
reconstruction pathways is resilient in precisely the sense
$\rho_W(r) \gg 1$ would formalize, and a root object with exactly one
surviving copy in one archive is the $\rho_W(r) = 1$ case this monograph's
oral-tradition discussion (Appendix~\ref{app:extended-examples}) already
worried about informally, without a name for it. This monograph does not
adopt $\rho_W(r)$ as its own machinery, since it has not been verified;
it is recorded here as a specific, checkable candidate for future work
rather than left as a vague gesture toward ``some measure or other.''
\end{remark}

\section{Where This Monograph's Apparatus Sits}
\label{sec:where-this-sits}


Five correspondences concerning Distinguishability Geometry's four
operations are worth stating explicitly, because they were not
recognized when the relevant chapters of this monograph were drafted, and
because they differ sharply in how tight they are. Four are proved
correspondences; one is an open, hedged analogy. Two further
correspondences, concerning evaluability and correction rather than the
four operations, were given already in \S\ref{sec:abt} and \S\ref{sec:goc}
and are not repeated here.

\begin{proposition}[Rendering Is Projection]
\label{prop:rendering-is-projection}
A rendering $r$ in the sense of \cref{def:rendering} is a projection in the
sense of \cref{def:dg-projection}, and the discard set $\discard{r}$ of
\cref{def:discard-set} corresponds to the fibers merged by that projection.
\end{proposition}

\begin{proofsketch}
Both are surjective maps from a source distinguishability structure to a
coarser one, defined by which pairs of source states become identified.
\cref{def:rendering}'s ``some distinctions separated in $X$ become identified
in $Y$'' is \cref{def:dg-projection}'s coarsening condition stated in
different words. What this monograph's apparatus does not compute, and
\cref{thm:projection-deficit-bound} does, is a quantitative lower bound on
how much is lost: $\Delta\delta(r) \ge H(X \mid Y)$ would give every discard
set in this book a numeric floor it currently lacks. This is a genuine,
available strengthening that has not been carried out.
\end{proofsketch}

\begin{proposition}[Augmentation Is Embedding]
\label{prop:augmentation-is-embedding}
Augmentation with an enumerable hidden-state set $H$
(\cref{prop:closed-sufficiency}) is a faithful embedding in the sense of
\cref{def:dg-embedding}, and \cref{prop:closed-sufficiency} is an instance of
\cref{thm:embedding-refinement}.
\end{proposition}

\begin{proofsketch}
The augmented rendering $r' = (r, H)$ is obtained from $r$ by an injective
map into a richer space that separates states $r$ identified, exactly
\cref{def:dg-embedding}'s shape; it is faithful precisely when $H$ exhausts
the task-relevant hidden state, which is \cref{prop:closed-sufficiency}'s
hypothesis. \cref{thm:embedding-refinement} then gives, for free and in full
generality, the conclusion \cref{prop:closed-sufficiency} proved by a bespoke
argument: a faithful embedding cannot increase the coding deficit, so
augmentation cannot make sufficiency worse. This monograph's proof of
\cref{prop:closed-sufficiency} is correct but is now known to be a special
case of a general theorem it did not cite.
\end{proofsketch}

\begin{proposition}[Relay Replacement Is Revision, Where the State Space Is Held Fixed]
\label{prop:relay-replacement-partial-revision}
Relay revision (\cref{def:relay-revision}) in which the underlying state
space $X$ is unchanged and only the descriptive ontology is expanded is an
instance of \cref{def:dg-revision}, and \cref{thm:revision-monotonicity}
guarantees the new theory's coding deficit is no worse than the old one's.
\end{proposition}

\begin{proofsketch}
When a superseded theory is replaced by one built on an expanded, but not
qualitatively different, evidential ontology -- more precise instruments
within the same conceptual framework, say -- this is exactly $T \subset T'$
on a fixed $X$, and \cref{thm:revision-monotonicity} applies directly. The
correspondence is deliberately not claimed in full generality. Where a
replacement changes what counts as a state at all -- a genuine paradigm
shift in the sense Chapter~\ref{ch:science}'s note on adjacent debates
explicitly declined to adjudicate -- the transition is not pure revision in
\cref{def:dg-revision}'s sense, since that definition holds $X$ fixed; such a
case would need to be analyzed as some composition of revision and
transport, and this monograph has not carried out that analysis.
\end{proofsketch}

\begin{proposition}[Paradigm-Shift Relay Replacement Is Repair-Repertoire Discontinuity]
\label{prop:relay-replacement-paradigm-shift}
Relay revision in which the underlying production process changes so
substantially that neither the old nor the new process's repertoire
contains the other corresponds to \cref{thm:paradigm-shift}, with the
substrate chain's producing repertoire at each era playing the role of
$\mathcal{R}_{\mathrm{old}}$ and $\mathcal{R}_{\mathrm{new}}$.
\end{proposition}

\begin{proofsketch}
This is exactly the case \cref{prop:relay-replacement-partial-revision}
left uncharacterized. \cref{thm:paradigm-shift} shows the resulting
restorable-distinction sets are incomparable rather than one being a
refinement of the other, which is a more decisive statement than
``analysis not carried out'': it says the two eras' restorable
distinctions are genuinely incommensurable in the technical sense of
neither set containing the other, with concrete lost and gained
distinctions on both sides. Chapter~\ref{ch:science}'s example of a
theory's evidential basis being enlarged rather than replaced falls under
\cref{prop:relay-replacement-partial-revision} instead; the paradigm-shift
case this proposition covers is the sharper one, where old and new
genuinely talk past each other on some distinctions in both directions.
\end{proofsketch}

\begin{remark}[Root Recovery as Persistence-Set Membership]
\label{rem:root-recovery-persistence-set}
\cref{prop:root-recovery} is naturally read as the claim that a specific
transformation -- the composite production process from $A_0$ to $A_n$ --
lies in the persistence set $P(d)$ (\cref{def:persistence-set}) for a
discarded distinction $d$, exactly when the reconstruction operator ``recompute
from the preserved $A_0$ under a revised policy'' is available in
$\mathcal{R}$. This reframes Appendix~\ref{app:open-questions}'s
root-preservation bootstrapping question sharply: rather than asking
qualitatively whether recovery is possible, one can ask what $\mu(P(d))$ is
-- what fraction of the relevant transformation family still admits
recovery -- and this monograph has not attempted that computation.
\end{remark}

\begin{remark}[Fidelity and Transport: An Open Correspondence]
\label{rem:fidelity-transport-open}
Chapter~\ref{ch:fidelity-continuum}'s fidelity $\fidelity{r}{\tau}$
(\cref{def:fidelity}) and \cref{thm:transport-stability}'s invariant-oscillation
bound both measure how well structure survives a transformation, but they
are different mathematical objects -- a cardinality ratio of admissible
continuation sets against a Lipschitz-style bound on a fiber-constant
invariant -- and this monograph has not shown one reduces to the other. This
is the weakest of the five correspondences in this section, offered as an
open question rather than a result, and recorded as such in
Appendix~\ref{app:open-questions}.
\end{remark}

\section{An Honest Scorecard}

Before this chapter, a reader could reasonably have concluded that this
monograph recognized only one of Distinguishability Geometry's four
primitive operations -- projection -- as load-bearing, with embedding,
revision, and transport absent. That conclusion would have been too
pessimistic and not quite accurate. \cref{prop:rendering-is-projection}
confirms projection was there from Chapter~\ref{ch:distinctions} onward.
\cref{prop:augmentation-is-embedding} shows embedding was present, unnamed,
in the augmentation mechanism of Chapters~\ref{ch:admissibility} and
\ref{ch:chess}. \cref{prop:relay-replacement-partial-revision} and
\cref{prop:relay-replacement-paradigm-shift} together show revision was
present, unnamed, in Chapter~\ref{ch:science} for both the case
Distinguishability Geometry's revision operation covers and the sharper
paradigm-shift case it does not, with Observability Is Restorability
supplying the missing half. Only transport remains without a demonstrated
instance in this monograph rather than merely an unnamed one:
\cref{rem:fidelity-transport-open}'s candidate is an analogy, not yet a
correspondence. Three operations out of four fully accounted for, and the
fourth genuinely open rather than merely unexamined, is a materially
different and more encouraging finding than one operation out of four --
and it is the finding this monograph is prepared to defend, rather than a
stronger claim it has not earned.

Two further findings sit outside the four-operations scorecard entirely.
\cref{prop:sufficiency-presupposes-evaluability} shows this monograph's
sufficiency apparatus was silently assuming an evaluability condition
Admissibility Before Truth makes explicit, discovered only once that
source was transcribed rather than paraphrased. \cref{prop:corruption-as-suppression}
shows this monograph's oldest and most-used distinction -- benign versus
corrupting loss, present since Chapter~\ref{ch:distinctions} -- is the
same phenomenon The Geometry of Correction names correction suppression,
recognized only by accident of the author happening to transcribe a
second, unrelated-seeming source. Neither finding would have been visible
from the four-operations framing alone, which is itself a reason to keep
checking new sources against this monograph's apparatus rather than
treating the reconciliation work in this chapter as finished.
